\exercisesection{\S~2}

\begin{enumerate}
\def\labelenumi{\arabic{enumi})}
\item
  The diagonal matrices of trace 0 form a Cartan subalgebra of \(\mathfrak{sl}(n,k)\) , except when \(n = 2\) and \(k\) is of characteristic 2.
\item
  Let \(e\) be the element \(\begin{pmatrix} 0 & 1 \\ 0 & 0 \end{pmatrix}\) of \(\mathfrak{sl}(2, \mathbf{C})\) . Show that \(\mathbf{C}e\) is a maximal nilpotent Lie subalgebra of \(\mathfrak{sl}(2, \mathbf{C})\) , but not a Cartan subalgebra of \(\mathfrak{sl}(2, \mathbf{C})\) .
\item
  Assume that \(k\) is of characteristic 0. Let \(\mathfrak{g}\) be a semi-simple Lie algebra. Let \(E\) be the set of commutative subalgebras of \(\mathfrak{g}\) all of whose elements are semi-simple in \(\mathfrak{g}\) . Then the Cartan subalgebras of \(\mathfrak{g}\) are the maximal elements of \(E\) . (Use Th. 2 and Prop. 10.)
\end{enumerate}

In particular, the union of the Cartan subalgebras of \(\mathfrak{g}\) is equal to the set of semi-simple elements of \(\mathfrak{g}\) .

\begin{enumerate}
\def\labelenumi{\arabic{enumi})}
\setcounter{enumi}{3}
\item
  Let \(\mathfrak{g}\) be a Lie algebra with a basis \((x,y,z)\) such that \([x,y] = y\) , \([x,z] = z\) , \([y,z] = 0\) . Let \(\mathfrak{a}\) be the ideal \(ky + kz\) of \(\mathfrak{g}\) . Then \(\mathrm{rk}(\mathfrak{a}) = 2\) and \(\mathrm{rk}(\mathfrak{g}) = 1\) .
\item
  Assume that \(k\) is of characteristic 0. Let \(\mathfrak{g}\) be a Lie algebra, \(\mathfrak{r}\) its radical, \(\mathfrak{h}\) a Cartan subalgebra of \(\mathfrak{g}\) . Show that
\end{enumerate}

\[
\mathfrak {r} = [ \mathfrak {g}, \mathfrak {r} ] + (\mathfrak {h} \cap \mathfrak {r}).
\]

(Observe that the image of \(\mathfrak{h}\) in \(\mathfrak{g}/[\mathfrak{g},\mathfrak{r}]\) contains the centre \(\mathfrak{r}/[\mathfrak{g},\mathfrak{r}]\) of \(\mathfrak{g}/[\mathfrak{g},\mathfrak{r}]\) .)

\begin{enumerate}
\def\labelenumi{\arabic{enumi})}
\setcounter{enumi}{5}
\item
  Let \(\mathfrak{g}\) be a Lie algebra, \(\mathfrak{h}\) a nilpotent subalgebra of \(\mathfrak{g}\) . If \(\mathfrak{g}^0 (\mathfrak{h})\) is nilpotent, \(\mathfrak{g}^0 (\mathfrak{h})\) is a Cartan subalgebra of \(\mathfrak{g}\) .
\item
  Let \(\mathfrak{s}\) be a Lie algebra, \(\mathfrak{a}\) a Cartan subalgebra of \(\mathfrak{s}\) and V an \(\mathfrak{s}\) -module. Let \(\mathfrak{g} = \mathfrak{s} \times V\) be the semi-direct product of \(\mathfrak{s}\) by V. Show that \(\mathfrak{a} \times V^0(\mathfrak{a})\) is a Cartan subalgebra of \(\mathfrak{g}\) .
\item
  Assume that \(k\) is of characteristic \(p > 0\) . Denote by \(\mathfrak{s}\) a Lie algebra with basis \(\{x, y\}\) such that \([x, y] = y\) . Let \(V\) be a \(k\) -vector space with basis \(\{e_i\}_{i \in \mathbf{Z}/p\mathbf{Z}}\) .
\end{enumerate}

\begin{enumerate}
\def\labelenumi{\alph{enumi})}
\item
  Show that V has a unique \(\mathfrak{s}\) -module structure such that \(xe_{i} = ie_{i}\) and \(ye_{i} = e_{i + 1}\) for all \(i\) . This \(\mathfrak{s}\) -module is simple.
\item
  Let \(\mathfrak{g} = \mathfrak{s} \times \mathrm{V}\) be the semi-direct product of \(\mathfrak{s}\) by \(\mathrm{V}\) . Show that \(\mathfrak{g}\) is a solvable algebra of rank 1 whose derived algebra is not nilpotent.
\item
  An element of \(\mathfrak{g}\) is regular if and only if its projection on \(\mathfrak{s}\) is of the form \(ax + by\) , with \(ab \neq 0\) .
\item
  We have \(\mathrm{V}^{0}(x+y)=0\) and \(\mathrm{V}^{0}(x)=ke_{0}\) . Deduce (cf.~Exerc. 6) that g has Cartan subalgebras of dimension 1 (for example that generated by \(x+y\) ) and Cartan subalgebras of dimension 2 (for example that generated by x and \(e_{0}\) ).
\end{enumerate}

\begin{enumerate}
\def\labelenumi{\arabic{enumi})}
\setcounter{enumi}{8}
\item
  Let g be a Lie algebra with a basis \((x,y)\) such that \([x,y]=y\) . Let k=ky, and \(\varphi:g\to g/k\) the canonical morphism. The element 0 of g/k is regular in g/k but is not the image under \(\varphi\) of a regular element of g.
\item
  Assume that \(k\) is infinite. Let \(\mathfrak{g}\) be a Lie algebra. Prove the equivalence of the following properties:
\end{enumerate}

\begin{enumerate}
\def\labelenumi{(\roman{enumi})}
\item
  \(\operatorname{rk}(\mathfrak{g}) = \dim (\mathfrak{g})\) .
\item
  \(\mathfrak{g}\) is nilpotent.
\item
  g has only finitely-many Cartan subalgebras of dimension rk(g).
\item
  g has only one Cartan subalgebra.
\end{enumerate}

\begin{enumerate}
\def\labelenumi{\arabic{enumi})}
\setcounter{enumi}{10}
\tightlist
\item
  Let \(\mathfrak{h}\) be a commutative Lie algebra \(\neq 0\) , \(P\) a finite subset of \(\mathfrak{h}^*\) containing 0. Show that there exists a Lie algebra \(\mathfrak{g}\) containing \(\mathfrak{h}\) as a Cartan subalgebra, and such that the set of weights of \(\mathfrak{h}\) in \(\mathfrak{g}\) is \(P\) . (Construct \(\mathfrak{g}\) as the semi-direct product of \(\mathfrak{h}\) by the \(\mathfrak{h}\) -module \(V\) which is the direct sum of the 1-dimensional modules corresponding to the elements of \(P - \{0\}\) , cf.~Exerc. 7.)
\end{enumerate}

An element \(x\) of \(\mathfrak{h}\) is such that \(\mathfrak{h} = \mathfrak{g}^0(x)\) if and only if \(x\) is not orthogonal to any element of \(\mathrm{P} - \{0\}\) .

\begin{enumerate}
\def\labelenumi{\arabic{enumi})}
\setcounter{enumi}{11}
\tightlist
\item
  Assume that \(k\) is finite. Construct an example of a Lie algebra \(\mathfrak{g}\) having a Cartan subalgebra \(\mathfrak{h}\) in which there exists no element \(x\) such that \(\mathfrak{h} = \mathfrak{g}^0(x)\) . (Use the preceding exercise, and take \(\mathrm{P} = \mathfrak{h}^*\) .)
\end{enumerate}

¶ 13) Assume that k is finite. Denote by \(k'\) an infinite extension of k. Let g be a Lie algebra over k. The rank of g, denoted by \(\operatorname{rk}(\mathfrak{g})\) , is the rank of the \(k'\) -Lie algebra \(g' = g \otimes_{k} k'\) ; an element of g is said to be regular if it is regular in \(g'\) ; these definitions do not depend on the choice of \(k'\) . Show that, if

\[
\operatorname{Card} (k) \geq \dim \mathfrak {g} - \operatorname{rk} (\mathfrak {g}),
\]

\(\mathfrak{g}\) contains a regular element (hence also a Cartan subalgebra).

(Use the following result: if \(a\) is a non-zero homogeneous element of \(k[\mathrm{X}_1, \ldots, \mathrm{X}_n]\) , and if \(\operatorname{Card}(k) \geq \deg(a)\) , there exists \(x \in k^n\) such that \(a(x) \neq 0\) .)

\begin{enumerate}
\def\labelenumi{\arabic{enumi})}
\setcounter{enumi}{13}
\tightlist
\item
  Assume that \(k\) is of characteristic zero. Let V be a finite dimensional \(k\) -vector space, \(\mathfrak{g}\) a Lie subalgebra of \(\mathfrak{gl}(V)\) , \(\mathfrak{h}\) a Cartan subalgebra of \(\mathfrak{g}\) and \(\mathfrak{n}_{\mathrm{V}}\) the largest ideal of nilpotency of the \(\mathfrak{g}\) -module V (Chap. I, §4, no. 3, Def. 2).
\end{enumerate}

Show that an element of h is nilpotent if and only if it belongs to \(n_{V}\) . (Reduce to the case where the g-module V is semi-simple, and use Cor. 3 of Th. 2.)

¶ 15) Assume that \(k\) is infinite. Let \(\mathfrak{g}\) be a Lie algebra, \(\mathcal{C}_{\infty}\mathfrak{g}\) the union of the ascending central series of \(\mathfrak{g}\) , and \(x\) an element of \(\mathfrak{g}\) . Prove the equivalence of the following properties:

\begin{enumerate}
\def\labelenumi{(\roman{enumi})}
\item
  \(x\) belongs to every Cartan subalgebra of \(\mathfrak{g}\) .
\item
  \(x\in \mathfrak{g}^0 (y)\) for all \(y\in \mathfrak{g}\) (i.e.~\(x\in \mathfrak{g}^0 (\mathfrak{g})\) ).
\item
  \(x \in C_{\infty}g.\)
\end{enumerate}

(The implications (iii) \(\Longrightarrow\) (ii) \(\Longrightarrow\) (i) are immediate. To prove that (i) \(\Longrightarrow\) (ii), remark that (i) is equivalent to saying that \(x\in \mathfrak{g}^0 (y)\) for every regular element \(y\) in \(\mathfrak{g}\) , and use the fact that the regular elements are dense in \(\mathfrak{g}\) in the Zariski topology. To prove that (ii) \(\Longrightarrow\) (iii), observe that \(\mathfrak{n} = \mathfrak{g}^0 (\mathfrak{g})\) is stable under \(\mathfrak{g}\) (§1, Exerc. 5) and apply Engel's theorem to the \(\mathfrak{g}\) -module \(\mathfrak{n}\) ; deduce that \(\mathfrak{n}\) is contained in \(\mathcal{L}_{\infty}\mathfrak{g}\) .)

¶ 16) Let g be a solvable complex Lie algebra, h a Cartan subalgebra of g, \(\mathfrak{g} = \bigoplus \mathfrak{g}^{\lambda}(\mathfrak{h})\) the corresponding decomposition of g into primary subspaces, with \(\mathfrak{g}^{0}(\mathfrak{h}) = \mathfrak{h}\) .

\begin{enumerate}
\def\labelenumi{\alph{enumi})}
\item
  Show that the restrictions to \(\mathfrak{h}\) of the linear forms called roots of \(\mathfrak{g}\) in Chap. III, §9, Exerc. 17 c) are the weights of \(\mathfrak{h}\) in \(\mathfrak{g}\) , i.e.~the \(\lambda\) such that \(\mathfrak{g}^{\lambda}(\mathfrak{h}) \neq 0\) ; deduce that such a \(\lambda\) vanishes on \(\mathfrak{h} \cap \mathcal{D}\mathfrak{g}\) .
\item
  Let \((x,y)\mapsto [x,y]'\) be the alternating bilinear map from \(\mathfrak{g}\times \mathfrak{g}\) to \(\mathfrak{g}\) with the following properties:
\end{enumerate}

\begin{enumerate}
\def\labelenumi{(\roman{enumi})}
\item
  If \(x \in \mathfrak{g}^{\lambda}(\mathfrak{h}), y \in \mathfrak{g}^{\mu}(\mathfrak{h})\) , with \(\lambda \neq 0, \mu \neq 0\) , then \([x, y]' = [x, y]\) ;
\item
  if \(x \in \mathfrak{g}^{0}(\mathfrak{h}), y \in \mathfrak{g}^{\mu}(\mathfrak{h})\) , then \([x, y]' = [x, y] - \mu(x)y\) .
\end{enumerate}

Show that this gives a new Lie algebra structure on \(\mathfrak{g}\) (use \(a\) ). Denote it by \(\mathfrak{g}'\) .

\begin{enumerate}
\def\labelenumi{\alph{enumi})}
\setcounter{enumi}{2}
\tightlist
\item
  Show that, if \(x \in \mathfrak{g}^{\lambda}(\mathfrak{h})\) , the map \(ad'x : y \mapsto [x, y]'\) is nilpotent. Deduce that \(g'\) is nilpotent (apply Exerc. 11 of Chap. I, §4 to the set E of \(ad'x\) , where x belongs to the union of the \(\mathfrak{g}^{\lambda}(\mathfrak{h})\) ).
\end{enumerate}

